%chapter2_image.tex
%----以下文字の大きさ，フォント，余白などを決定する設定----
\documentclass[a4paper,10pt]{jsarticle}
\setlength{\topmargin}{-17truemm}
\setlength{\oddsidemargin}{-0.4truemm}
\setlength{\textwidth}{160truemm}
\setlength{\textheight}{247truemm}
%-------------------------------------------------

%\usepackage{}：Latexで拡張機能を使用するための宣言-------------
\usepackage[dvipdfmx]{graphicx}%図を埋め込むためのパッケージ
\usepackage{float}%その場に表・図を表示させるパッケージ
%------------------------------------------------------------------

\title{画像を表示する}%タイトル
\author{ER18023 金田 康志}%著者
\date{2020年12月27日}%日付

\begin{document}%本文の始まり
\maketitle%上記で設定したtitle,author,dateを表示
\section{画像の表示について}
chapter2\_image.texとchapter2\_image.pdfを照らし合わせて読んでください．\\

usepackageでdvipdfmx, graphicx, floatを指定すると画像を埋め込むことが可能です．\\

図は，図\ref{fig:label}のように参照することが可能です．基本的に，図を記述する場合，必ず1度以上参照する必要があります．
また，一度コンパイルしただけだと，参照した時に「?」と表示されるので「pLaTeX(ptex2pdf)」→「pLaTeX(ptex2pdf)」と2回コンパイルを実行する必要があります．\\
※ VScodeでの環境構築が完了している場合，「ctrl+s」で保存するだけでコンパイルが完了します．\\

以下に挿入図を図\ref{fig:label}に示す
%---------------図の表示方法--------------------
\begin{figure}[h]
    \begin{center}
        \includegraphics[width=100mm]{sample1.png}
    \end{center}
    \caption{挿入図}
    \label{fig:label}
\end{figure}
%------------------------------------------

%---------------以下課題-------------------------------
\section{課題}
\noindent 本chapterは必須の課題のみになっています．\\ 
必須1：authorを自分の名前，dateを提出日に変更してください．\\
必須2：図を追加してください（図は何でも可）．また，refを使用して「図を図\ref{fig:label}に示す」の様に参照してください．\\

\noindent 編集が完了したら，コンパイルしてください．\\
VScodeで編集している場合，「ctrl+s」で保存＆コンパイル→「ctrl+alt+V」で出力されるpdfをプレビューできるので変更点が反映されていることを確認してください．
%--------------------------------------------------------
\end{document}

%以下，今回登場した機能の説明
%----------------------図----------------------------------------------
%\begin{figure}[画像の表示位置]                    　　　% 表の表示位置の例→h：その場，b：ページ下部，t：ページ上部，p：他のページ
%    \begin{center}                               　　　% \begin{center}から\end{center}の間までが中心に表示される
%        \includegraphics[width=横幅]{画像のパス}       % 画像までのパスは相対パス
%    \end{center}
%    \caption{図のタイトル}                                                
%    \label{ラベル名}　　　　　　　　　　　　　　　　　　　%図の場合は\label{fig:ラベル名}で定義することが多い．ラベルは\ref{ラベル名}で参照可能．
%\end{figure}
%---------------------------------------------------------------------